\section{TPU/Pallas Kernel Optimization}

\smallskip\noindent\textit{Study Guide companion: Study Guide Section 12 (Hardware Deep Dive) covers the interview-ready version of this material.}
\smallskip

The RLVR training pipeline is ``held together by duct tape.'' Every percentage point of hardware utilization matters at scale. Simon Boehm's blog post walking through CUDA matmul optimization from 1.3\% to 93.7\% of cuBLAS remains the reference for GPU kernel pedagogy. No equivalent guide exists for TPUs and Pallas. This chapter follows the same methodology: start naive, profile, fix one thing, measure, repeat.

\subsection{Conceptual Foundation}

\begin{thinkfirst}{GPU vs.\ TPU mental models}
The biggest conceptual shift: on GPUs you think about \textbf{threads and memory access patterns}. On TPUs you think about \textbf{pipeline efficiency and keeping the systolic array fed}.

\begin{enumerate}
  \item \textbf{The GPU model:} Thousands of lightweight threads execute the same instruction on different data (SIMT). The bottleneck is usually memory bandwidth --- getting data to the threads fast enough. Key optimizations: coalescing memory accesses, using shared memory, avoiding bank conflicts.

  \item \textbf{The TPU model:} A systolic array (128$\times$128 or 256$\times$256 multiply-accumulate units) flows data through in a deterministic pattern. There's no warp scheduling randomness. The bottleneck is keeping data flowing through the array without bubbles. Key optimizations: tiling to match MXU dimensions, pipelining HBM-to-VMEM transfers, maximizing arithmetic intensity.

  \item \textbf{Why this matters for RLVR:} During RLVR training, the pipeline alternates between generation (autoregressive, memory-bandwidth-bound) and training (matmul-heavy, compute-bound). RLVR runs are particularly long because verification-based rewards require generating complete solutions and checking them end-to-end --- the generation phase is therefore even longer per update than in simpler RL setups. The optimization strategy is different for each phase. On a TPU, the generation phase underutilizes the MXU.

  Now answer: if RLVR training spends 60\% of time on generation and 40\% on gradient updates, which phase should you optimize first? What's the theoretical maximum speedup from optimizing just that phase? (Hint: Amdahl's law.)
\end{enumerate}
\end{thinkfirst}

\begin{thinkfirst}{Memory hierarchy reasoning}
TPU memory hierarchy: HBM (off-chip, large, slow) $\to$ VMEM (on-chip SRAM, small, fast) $\to$ VREGs (registers).

\begin{enumerate}
  \item A TPU v5e has 16GB HBM at 820 GB/s bandwidth and $\sim$30 MB VMEM. You want to multiply two 4096$\times$4096 bf16 matrices. Each matrix is 32 MB --- neither fits in VMEM. How do you tile this computation so each tile fits in VMEM while maximizing MXU utilization?

  \item The MXU does 128$\times$128 matmuls per cycle. If your tile is 64$\times$64, what fraction of the MXU is utilized? What about 256$\times$256? Work out the utilization for each and explain the tradeoff between tile size and MXU efficiency.

  \item Pallas uses ``BlockSpec'' to define tiles and ``Buffered(buffer\_count=2)'' to double-buffer HBM$\to$VMEM transfers. Draw the pipeline on paper: while tile $N$ is being computed on the MXU, tile $N+1$ is being loaded into the second VMEM buffer via DMA. What determines whether you're compute-bound or memory-bound in this pipelined setup?
\end{enumerate}
\end{thinkfirst}

\begin{thinkfirst}{Arithmetic intensity and the roofline model}
Arithmetic intensity = FLOPs / bytes transferred. For a matrix multiply: $2mnk$ FLOPs, $(mn + mk + nk) \times \text{bytes\_per\_element}$ transferred. For square matrices: intensity $\approx n/3$ (for bf16).

\begin{enumerate}
  \item A TPU v5e has $\sim$200 TFLOP/s compute and 820 GB/s memory bandwidth. At what matrix size does a bf16 matmul become compute-bound rather than memory-bound? Use the roofline model: compute the crossover arithmetic intensity, then solve for $n$.

  \item RLVR training involves many small matmuls (e.g., reward model forward passes on individual completions). Are these compute-bound or memory-bound? What does this mean for optimization strategy?

  \item The autoregressive generation phase produces one token at a time. Each step involves a matmul of shape $[1, d] \times [d, V]$ where $V$ is vocabulary size. Calculate the arithmetic intensity. Is this compute-bound or memory-bound? (This is why generation is the bottleneck --- it's always memory-bandwidth-bound.)
\end{enumerate}
\end{thinkfirst}

\subsection{Hands-On Exercises}

\begin{exercise}{Pallas matmul: naive to optimized (10--15 hours)}
Requires: TPU access (Google Cloud TPU Research program or Colab TPU).

Write a Pallas matmul kernel and optimize it step by step, measuring TFLOP/s at each stage:

\textbf{Kernel 0: JAX baseline.} \texttt{jnp.dot(A, B)}. This is XLA's auto-tuned kernel. Record as your ceiling.

\textbf{Kernel 1: Naive Pallas.} Single block, no tiling. Process the entire matrix in one kernel invocation.
\begin{lstlisting}
def naive_matmul_kernel(x_ref, y_ref, z_ref):
    z_ref[...] = x_ref[...] @ y_ref[...]

result = pl.pallas_call(
    naive_matmul_kernel,
    out_shape=jax.ShapeDtypeStruct((m, n), jnp.float32),
    grid=(1,),
    in_specs=[pl.BlockSpec((m, k), lambda i: (0, 0)),
              pl.BlockSpec((k, n), lambda i: (0, 0))],
    out_specs=pl.BlockSpec((m, n), lambda i: (0, 0)),
)(x, y)
\end{lstlisting}

\textbf{What to measure:} TFLOP/s. This will be terrible because the matrices probably don't fit in VMEM.

\textbf{Kernel 2: Block tiling.} Add a grid and proper BlockSpecs. Tile along the reduction dimension.

\textbf{Kernel 3: FP32 accumulator.} Use bf16 inputs but accumulate in FP32 (use \texttt{scratch\_shapes} for the accumulator in VMEM).

\textbf{Kernel 4: Pipeline the DMA.} Add \texttt{Buffered(buffer\_count=2)} to overlap DMA with compute.

\textbf{Kernel 5: Block size sweep.} Try $bm, bn, bk \in \{64, 128, 256, 512\}$. Plot TFLOP/s vs.\ block size. Why does 128 usually win? (It matches the MXU dimension.)

\textbf{Kernel 6: Megacore.} Set \texttt{dimension\_semantics=("parallel", "parallel", "arbitrary")} to use both TensorCores. Expected: $\sim$2$\times$ on supported chips.

\textbf{At each step:} Record TFLOP/s and \% of JAX baseline. Each optimization should produce a measurable improvement --- if it doesn't, diagnose why before moving on. By Kernel 4--5, you should be at $>50\%$ of the JAX baseline. If you're below 20\% after Kernel 3, something fundamental is wrong with your tiling.
\end{exercise}

\begin{exercise}{Extended Pallas Kernels (6--8 hours)}
Continue directly from Kernel 6. These four kernels push into production-quality territory: fused ops, mixed precision, chip-specific tuning, and distributed execution.

\textbf{Kernel 7: Fused transpose + activation.} Write a single Pallas kernel that applies a matrix transpose followed by an element-wise activation (GELU or SiLU) without writing the intermediate transposed result back to HBM.

\begin{itemize}
  \item Load a tile from HBM, transpose it in VMEM registers, apply the activation, write the fused result.
  \item \textbf{Measure:} TFLOP/s improvement vs.\ two separate kernels (transpose then activation). What fraction of the speedup comes from eliminating the intermediate HBM write?
  \item \textbf{Think:} GELU requires $\tanh$ --- how does the TPU handle transcendental functions? Is it approximated? Does SiLU (which avoids $\tanh$) outperform GELU on TPUs specifically?
\end{itemize}

\textbf{Kernel 8: Custom precision paths.} bf16 input, fp32 accumulate, bf16 output.

\begin{itemize}
  \item Use \texttt{scratch\_shapes} to allocate an fp32 accumulator in VMEM. Accept bf16 inputs via \texttt{in\_specs}, cast to fp32 on load, accumulate, cast back to bf16 on store.
  \item \textbf{Measure:} Compare gradient norm variance across 100 training steps vs.\ pure bf16 accumulation. Numerical stability improvement should be visible if your learning rate is high or your model is deep.
  \item \textbf{Think:} Pure bf16 accumulation introduces quantization error on every add. At what model size or learning rate does this become practically significant? Cross-check against JAX's default XLA behavior --- does XLA already do this internally?
\end{itemize}

\textbf{Kernel 9: TPUv6e-specific optimization (if accessible).} TPUv6e has a larger MXU and more VMEM than v5e. Re-tune block sizes for v6e architecture.

\begin{itemize}
  \item Re-run your Kernel 5 block size sweep on v6e. Does the optimal $(bm, bn, bk)$ change?
  \item If v6e is not accessible, use the TPU Research Cloud application to request access (this is worth doing regardless --- document the application process as part of your writeup).
  \item \textbf{Think:} If the optimal block size scales with MXU dimensions, what block sizes would you predict for v6e before running the sweep? Make the prediction first, then measure.
\end{itemize}

\textbf{Kernel 10: Distributed matmul with RDMA.} Promote the ring all-reduce exercise below into a full kernel step: implement a distributed matmul where each device computes a shard and communicates partial results.

\begin{itemize}
  \item Partition matrix $A$ along rows across $N$ devices. Each device computes its shard of $C = A_{\text{shard}} \cdot B$. Use ring RDMA (see the distributed exercise below) to accumulate partial sums.
  \item \textbf{Measure:} Effective TFLOP/s per device as $N$ scales from 1 to 8. Does it scale linearly? Where does communication overhead dominate?
  \item \textbf{Think:} This is a simplified version of tensor parallelism. What would need to change to make this production-viable in a real RLVR training loop?
\end{itemize}
\end{exercise}

\begin{exercise}{Profile the bottleneck in RLVR training (4--6 hours)}
Profile an actual RLVR training step (use TRL's \texttt{GRPOTrainer} or your nanoRLVR from Chapter 8).

\begin{enumerate}
  \item Use \texttt{jax.profiler} or PyTorch profiler to trace one full GRPO update step.
  \item Break down time: what fraction is generation (sampling completions)? What fraction is forward pass (reward model or verifier)? What fraction is backward pass (policy gradient)?
  \item Identify the actual bottleneck. Is it what you expected? Note: because RLVR requires complete solution generation for verification, generation typically dominates even more than in RLHF-style training.
  \item For the bottleneck operation: is it compute-bound or memory-bound? Calculate the arithmetic intensity to verify.
  \item Propose \textit{one} optimization for the bottleneck. Implement it if you can. Measure the speedup.
  \item \textbf{Write a 1-paragraph memo:} You have one week of engineering time. Based on your profiling, which optimization would you prioritize for the RLVR training pipeline, and why? What's the expected speedup vs.\ engineering cost? This is the judgment call a senior RE makes daily.
\end{enumerate}

\end{exercise}

The point of this exercise: \textbf{measure before optimizing}. If you assumed the bottleneck before profiling and got it wrong, that's the lesson. Diagnosing the real bottleneck rather than the assumed one is exactly what ``attacking the problem directly'' looks like.

\begin{exercise}{Distributed Pallas: ring all-reduce (8--12 hours)}
Requires: multi-chip TPU access (TPU v4-8 or larger pod slice).

Implement a ring all-reduce in Pallas using RDMA primitives:

\begin{enumerate}
  \item Each device holds a shard of the gradient tensor.
  \item In $N$ steps (where $N$ = number of devices), each device sends its shard to the right neighbor and accumulates the incoming shard.
  \item After $N$ steps, every device has the complete sum.
\end{enumerate}

\textbf{Key Pallas primitives:}
\begin{lstlisting}
copy = pltpu.make_async_remote_copy(
    src_ref=local_buf,
    dst_ref=remote_buf,
    send_sem=send_sem,
    recv_sem=recv_sem,
    device_id=right_neighbor
)
copy.start()    # Non-blocking
# ... compute on local data ...
copy.wait_recv()  # Block until received
\end{lstlisting}

\textbf{Questions to answer:}
\begin{enumerate}[label=\alph*)]
  \item Derive the theoretical bandwidth utilization of ring all-reduce as a function of $N$ (number of devices). Does it scale well?
  \item Measure actual vs.\ theoretical bandwidth. Where's the gap?
  \item Can you overlap the all-reduce communication with local computation? (Hint: start the DMA, compute on local data, then wait for completion.)
\end{enumerate}
\end{exercise}

\begin{thinkfirst}{TPUv7 Ironwood: Inference-First Chip for Training?}
TPUv7 Ironwood is Google's latest TPU, designed as an inference-first chip: 192GB HBM3E, optimized for large-scale serving workloads. Anthropic has access to 1M Ironwoods and runs RLVR training on TPUs. This creates a tension worth reasoning through carefully.

\begin{enumerate}
  \item \textbf{The roofline question.} Ironwood has 192GB HBM3E --- the same capacity as MI300X. But memory \textit{bandwidth} and peak \textit{compute} may be balanced differently from v5e. Recall the roofline crossover: ridge point = peak compute / memory bandwidth. If Ironwood has higher bandwidth but the same or lower peak compute than v5e, what happens to the ridge point? Does this make it easier or harder to stay in the compute-bound regime for training matmuls?

  \item \textbf{RLVR's dual personality.} RLVR training interleaves two very different workloads: generation (autoregressive, memory-bandwidth-bound) and gradient updates (dense matmul, compute-bound). If Ironwood is optimized for the inference-like generation phase, does it actually \textit{help} RLVR more than a training-optimized chip would? Sketch the utilization profile across a single RLVR step on Ironwood vs.\ v5e.

  \item \textbf{Argue both sides.} The inference-first design may feature higher memory bandwidth but lower peak compute. For RLVR where generation dominates wall time:
  \begin{itemize}
    \item \textbf{Argument FOR Ironwood being the right chip:} Generation is memory-bandwidth-bound; higher bandwidth directly reduces wall time in the dominant phase. RLVR's bottleneck is where Ironwood is strongest.
    \item \textbf{Argument AGAINST:} The gradient update phase, even if shorter, still gates convergence. If Ironwood is significantly weaker on dense matmuls, the training phase takes longer per batch, limiting how fast you can improve the policy. A training-optimized chip that keeps the ratio of generation-to-update time better balanced might converge in fewer wall-clock hours overall.
  \end{itemize}
  Which argument do you find more convincing? What empirical measurement would settle the question?

  \item \textbf{The 1M chip question.} Anthropic reportedly has 1M Ironwoods. At that scale, even a 5\% improvement in RLVR throughput has enormous leverage. What single kernel optimization from this chapter --- if ported to Ironwood --- would have the highest expected impact? (You may not have Ironwood access; reason from first principles about the architecture.)
\end{enumerate}
\end{thinkfirst}

\begin{portfolio}{TPU Optimization Guide: ``The Pallas Companion''}
Write the Simon Boehm blog post but for TPUs. An iterative, pedagogical guide to writing high-performance Pallas kernels, with profiling at every step.

\textbf{Structure (10 kernels):}
\begin{enumerate}
  \item JAX baseline (XLA auto-tuned)
  \item Naive Pallas (single block)
  \item Block tiling with grid
  \item FP32 accumulator in VMEM scratch
  \item Double-buffered pipelining
  \item Block size autotuning
  \item Megacore parallelism
  \item Fused operations (transpose, activation)
  \item Custom precision (bf16 input, fp32 accumulate)
  \item Distributed matmul with RDMA
\end{enumerate}

\textbf{What makes this portfolio-quality:}
\begin{itemize}
  \item Performance table at each step (\% of peak, TFLOP/s)
  \item Roofline model analysis showing compute vs.\ memory bound transitions
  \item Diagrams of memory hierarchy, pipeline stages, and tiling patterns
  \item Honest documentation of what didn't work and why
  \item Comparison of TPU vs.\ GPU optimization strategies
\end{itemize}

\textbf{The monkey:} ``Can you teach the optimization journey clearly enough that a reader with CUDA experience but no TPU experience reaches $>$50\% of peak on their first try?'' Writing the kernels is the pedestal. Writing the \textit{explanation} that makes the reader understand \textit{why} each optimization works is the monkey.

\textbf{Why this fills a gap:} Boehm's CUDA guide is ``the reference'' for GPU optimization. No equivalent exists for TPUs. The official Pallas docs are reference material, not pedagogy.

\textbf{Hardware note on Ironwood:} The results in this guide are developed on TPU v5e (and v6e if accessible via TRC). For Ironwood-specific results, apply to the Google TPU Research Cloud program and request v7 access when available. The findings on v5e extrapolate \textit{architecturally} but not \textit{quantitatively} to Ironwood --- the optimal block sizes, roofline ridge points, and bandwidth utilization numbers will differ due to the inference-first memory architecture. Documenting the delta between v5e and Ironwood predictions vs.\ measurements is itself a valuable contribution.

\textbf{Verification:} Find one person with CUDA experience and no TPU experience. Have them follow your guide. If they reach 50\% of peak TFLOP/s, you've succeeded.
\end{portfolio}

\newpage
